\documentclass[11pt]{article}
\input{style-pc}
\usepackage{array}

\begin{document}

\entetesujet{Sujet bac 2015 -- Physique-Chimie -- Série D}

% ============================ CHIMIE ============================
\matiere{CHIMIE}{8 points}

\partie{A : vérification des connaissances}

\noindent\textbf{1. Questions à choix multiples.} Choisis la bonne réponse.
\begin{enumerate}[label=\alph*.]
  \item Le pH d'une solution aqueuse d'une monobase forte de concentration molaire $C$ est :
  \begin{enumerate}[label=a\textsubscript{\arabic*})]
    \item $14 - \log C$
    \item $-\log C$
    \item $14 + \log C$
  \end{enumerate}
  \item Un oxydant est une espèce chimique capable :
  \begin{enumerate}[label=b\textsubscript{\arabic*})]
    \item de céder des électrons.
    \item de capter des électrons.
  \end{enumerate}
  \item Lorsqu'on dilue une solution acide :
  \begin{enumerate}[label=c\textsubscript{\arabic*})]
    \item le pH diminue ;
    \item le pH augmente ;
    \item le pH ne varie pas.
  \end{enumerate}
\end{enumerate}

\medskip
\noindent\textbf{2. Question à alternative vrai ou faux.} Réponds par vrai ou faux.
\begin{enumerate}[label=\alph*.]
  \item Le catalyseur permet de réduire la durée d'une réaction.
  \item Le temps de demi-réaction d'une réaction d'ordre 2 est $\dfrac{\ln 2}{K}$.
  \item L'hydrolyse d'un ester est une réaction limitée.
\end{enumerate}

\medskip
\noindent\textbf{3. Réarrangement.} La phrase suivante a été écrite en désordre : ordonne-la.

\textit{L'activité / par seconde / radioactive / est le nombre / d'une source / de désintégration}

\partie{B : application des connaissances}
La réaction de décomposition de l'azométhane $\ch{CH_3N_2CH_3}$ suivant l'équation :
\[ \ch{CH_3N_2CH_{3(g)} \longrightarrow CH_3CH_{3(g)} + N_{2(g)}} \]
est une réaction d'ordre un. Sachant que la constante de vitesse est $K = 4\cdot10^{-4}$ s$^{-1}$ et que la concentration initiale est $C_0 = 0{,}604$ mol$\cdot$L$^{-1}$ :

\begin{enumerate}
  \item Écris la loi de vitesse de cette réaction.
  \item Calcule :
  \begin{enumerate}[label=\alph*.]
    \item Le temps de demi-réaction.
    \item Le temps nécessaire à la disparition de 75 \% de la concentration initiale du réactif.
  \end{enumerate}
  \item \begin{enumerate}[label=\alph*.]
    \item Détermine la concentration de l'azométhane à l'instant $t = 10$ min.
    \item Déduis la vitesse de la réaction à cet instant.
  \end{enumerate}
\end{enumerate}

% ============================ PHYSIQUE ============================
\matiere{PHYSIQUE}{12 points}

\partie{A : vérification des connaissances}

\noindent\textbf{1. Texte à trous.} Complète la phrase par les mots suivants : \textit{tangent ; trajectoire ; uniforme ; normal}.

\textit{Dans un mouvement circulaire $\cdots\cdots$, le vecteur accélération est $\cdots\cdots$ tandis que le vecteur vitesse est $\cdots\cdots$ à la $\cdots\cdots$}

\medskip
\noindent\textbf{2. Appariement.} Relie l'élément-question de la colonne A à l'élément-réponse de la colonne B qui lui correspond. Exemple : $a_5 = b_7$.

\begin{center}\renewcommand{\arraystretch}{1.8}
\begin{tabular}{|p{0.55\linewidth}|p{0.34\linewidth}|}
\hline
\textbf{Colonne A} & \textbf{Colonne B} \\
\hline
$a_1$ : impédance du circuit $R, L$ & $b_1$ : $Z = \sqrt{R^2 + \left(L\omega - \dfrac{1}{C\omega}\right)^2}$ \\
\hline
$a_2$ : différence de marche d'un point d'amplitude nulle & $b_2$ : $Z = \sqrt{R^2 + (L\omega)^2}$ \\
\hline
$a_3$ : impédance du circuit $R, L, C$ en série & $b_3$ : $d_2 - d_1 = k\lambda$ \\
\hline
$a_4$ : différence de marche d'un point d'amplitude maximale & $b_4$ : $d_2 - d_1 = (2k + 1)\dfrac{\lambda}{2}$ \\
\hline
\end{tabular}
\end{center}

\noindent\textbf{3. Questions à choix multiples.} Choisis la bonne réponse. Exemple : $f = f\cdot3$.
\begin{enumerate}[label=\alph*.]
  \item Dans un mouvement circulaire uniformément varié, l'accélération angulaire est :
  \begin{enumerate}[label=a\textsubscript{\arabic*})]
    \item nulle
    \item constante
    \item variable
  \end{enumerate}
  \item Le mouvement d'un satellite autour de la Terre est étudié dans un référentiel :
  \begin{enumerate}[label=b\textsubscript{\arabic*})]
    \item terrestre supposé galiléen
    \item géocentrique supposé galiléen
    \item lié au centre du Soleil
  \end{enumerate}
  \item La célérité $V$ des ondes qui se propagent sur une corde de masse linéique $\mu$, tenue grâce à une force d'intensité $F$ est :
  \begin{enumerate}[label=c\textsubscript{\arabic*})]
    \item $V = \sqrt{\dfrac{\mu}{F}}$
    \item $V = \sqrt{F\mu}$
    \item $V = \sqrt{\dfrac{F}{\mu}}$
  \end{enumerate}
  \item L'énergie cinétique d'un système animé d'un mouvement de rotation est donnée par l'expression :
  \begin{enumerate}[label=d\textsubscript{\arabic*})]
    \item $\dfrac{1}{2}mV^2$
    \item $\dfrac{1}{2}J\dot{\theta}^2$
    \item $\dfrac{1}{2}\dfrac{J}{\dot{\theta}^2}$
  \end{enumerate}
\end{enumerate}

\partie{B : application des connaissances}
Un camion de masse totale $M = 2{,}4$ tonnes grimpe une côte rectiligne $AB$, inclinée d'un angle $\alpha = 30^\circ$ par rapport à l'horizontale. Partant du repos de $A$, il accélère uniformément tel qu'il atteint la vitesse de 18 km/h en 10 secondes. Les forces de frottement sur ce trajet sont équivalentes à une force unique $\vec{f}$ parallèle à la ligne de plus grande pente dont l'intensité est : $f = 400$ N. Calcule :

\begin{enumerate}[label=\alph*.]
  \item L'accélération du mouvement du véhicule.
  \item L'intensité $F$ de la force motrice exercée par le moteur du camion.
  \item La vitesse du véhicule au sommet $B$ de la côte sachant que $AB = 196$ m.
  \item L'énergie mécanique $E_m$ du système (camion-Terre) au sommet $B$ de la côte. On prendra pour niveau de référence de l'énergie potentielle le plan horizontal passant par $A$.
\end{enumerate}

\partie{C : résolution d'un problème}
On désire connaître la longueur d'onde $\lambda$ d'une radiation lumineuse en exploitant les résultats d'une expérience portant sur l'effet photoélectrique. On dispose d'une cellule photoélectrique dont la cathode photo-émissive est caractérisée par un seuil photoélectrique correspondant à la longueur d'onde $\lambda_0 = 0{,}684\cdot10^{-6}$ m. On l'éclaire par la radiation de longueur d'onde $\lambda < \lambda_0$. On constate que, pour une différence de potentiel entre l'anode et la cathode égale à 45 V, les électrons émis arrivent sur l'anode avec une vitesse $v_A = 4\cdot10^6$ m/s.

\begin{enumerate}
  \item Détermine :
  \begin{enumerate}[label=\alph*.]
    \item L'énergie d'extraction $W_0$ d'un électron de la cathode.
    \item L'énergie cinétique d'un électron arrivant sur l'anode.
    \item L'énergie cinétique maximale d'un électron émis à la cathode.
  \end{enumerate}
  \item Calcule l'énergie $W$ d'un photon incident.
  \item Déduis-en la valeur de la longueur d'onde de cette radiation.
\end{enumerate}
\noindent Données : $m_e = 9{,}1\cdot10^{-31}$ kg ; $c = 3\cdot10^8$ m/s ; $h = 6{,}62\cdot10^{-34}$ J$\cdot$s ; $e = 1{,}6\cdot10^{-19}$ C.

\end{document}
